%%% BEGIN fig:dis-euclid-segment
\begin{figure}[htbp]
\centering
\begin{tikzpicture}[fd,
  fb bcut/.style={draw,line width=0.9pt,decorate,decoration={zigzag,segment length=4.2pt,amplitude=1.6pt}},
  fb seg/.style={draw=blue!65!black,line width=2.2pt,line cap=butt},
  fb lab/.style={font=\footnotesize},
  fb zoom/.style={draw=gray!55!black,line width=0.45pt}]
  \def\fbu{1.75}   % cm per unit of omega
  % ---------------- omega plane ----------------
  \fill[blue!5] (0,0) circle[radius=\fbu];
  \draw[gray!65!black,line width=0.7pt,dash pattern=on 1.2pt off 1.8pt] (0,0) circle[radius=\fbu];
  \draw[lcaxis] (-2.45*\fbu,0) -- (2.55*\fbu,0) node[above left=3pt and -2pt,font=\small] {$\Real\omega$};
  \draw[lcaxis] (0,-1.5*\fbu) -- (0,1.75*\fbu) node[left,font=\small] {$\Imag\omega$};
  \draw[fb bcut] (\fbu,0) -- (2.4*\fbu,0);
  \draw[fb bcut] (-\fbu,0) -- (-2.4*\fbu,0);
  \fill (\fbu,0) circle[radius=1.6pt]; \fill (-\fbu,0) circle[radius=1.6pt];
  \node[anchor=north west,font=\small,inner sep=1.5pt] at (\fbu+0.02,-0.08) {$1$};
  \node[anchor=north east,font=\small,inner sep=1.5pt] at (-\fbu-0.02,-0.08) {$-1$};
  \node[fb lab,anchor=north,align=center] at (1.9*\fbu,-0.16) {$s$ 道割线\\（物理区 $0<x\le1$）};
  \node[fb lab,anchor=north] at (-1.85*\fbu,-0.16) {$u$ 道割线};
  \node[fb lab,text=gray!35!black,anchor=north east] at (-0.72*\fbu,-0.72*\fbu) {$|\omega|=1$};
  \node[fb lab,align=left,text=blue!45!black,anchor=south east] (AN) at (-1.05,1.82) {单位圆盘 $|\omega|<1$ 内\\$T_2$ 解析（\cref{prop:dis-dispersion-np}）};
  \draw[blue!45!black,line width=0.4pt] ($(AN.south east)+(-0.25,0.02)$) -- (-0.72,0.9);
  % the Euclidean segment I_Q (drawn to scale for sqrt(3)M/Q = 0.08)
  \draw[fb seg] (0,-0.08*\fbu) -- (0,0.08*\fbu);
  % magnifier: small circle at the origin, external tangents to the inset
  \coordinate (Z) at (5.6,3.4);
  \draw[fb zoom] (0,0) circle[radius=0.49];
  \draw[fb zoom] (-0.316,0.375) -- (4.633,4.547);
  \draw[fb zoom] (0.187,-0.453) -- (6.171,2.013);
  \begin{scope}[shift={(Z)}]
    \begin{scope}
      \clip (0,0) circle[radius=1.5];
      \fill[blue!5] (-1.6,-1.6) rectangle (1.6,1.6);
      % the origin of omega sits at (-0.5,0) inside the inset
      \draw[gray!60!black,line width=0.5pt] (-1.6,0) -- (1.6,0);
      \draw[gray!60!black,line width=0.5pt] (-0.5,-1.6) -- (-0.5,1.6);
    \end{scope}
    \draw[fb zoom,line width=0.7pt] (0,0) circle[radius=1.5];
    \draw[fb seg] (-0.5,-0.86) -- (-0.5,0.86);
    \draw[blue!65!black,line width=0.6pt] (-0.61,0.86) -- (-0.39,0.86) (-0.61,-0.86) -- (-0.39,-0.86);
    \node[font=\scriptsize,anchor=west,inner sep=2pt] at (-0.4,0.86) {$\ii\sqrt3\,M/Q$};
    \node[font=\scriptsize,anchor=west,inner sep=2pt] at (-0.4,-0.86) {$-\ii\sqrt3\,M/Q$};
    \node[font=\small,text=blue!55!black,anchor=west] at (-0.38,0.36) {$\mathcal I_Q$};
    \draw[red!60!black,line width=0.7pt,-{Stealth[length=4pt,width=3.4pt]}] (-0.88,1.16) -- (-0.88,0.6);
    \draw[red!60!black,line width=0.7pt,-{Stealth[length=4pt,width=3.4pt]}] (-0.88,-1.16) -- (-0.88,-0.6);
  \end{scope}
  \node[font=\footnotesize,anchor=west,align=left] at ($(Z)+(1.62,0)$) {{\color{gray!25!black}原点附近放大}\\[3pt]{\color{red!55!black}$Q\to\infty$：}\\{\color{red!55!black}线段缩向原点}};
\end{tikzpicture}
\caption{前向 Compton 振幅 $T_2$ 在复 $\omega=1/x$ 平面上的欧氏点。靶静止系中，欧氏动量 $q=(\ii q^4,\mathbf q)$（$|\mathbf q|\ge Q/2$）给出 $\omega=2\ii Mq^4/Q^2$，
恰好取遍虚轴上的短线段 $\mathcal I_Q=\{\ii y:\,|y|\le\sqrt3M/Q\}$（蓝色粗线；右上为原点附近的放大），$Q\to\infty$ 时它缩向原点（VF-REF-1）。
$T_2$ 在单位圆盘 $|\omega|<1$ 内解析，割线 $(-\infty,-1]$ 与 $[1,\infty)$ 分别为 $u$ 道与 $s$ 道（物理区）的割线。
欧氏 OPE 只在长度为 $\order{M/Q}$ 的线段 $\mathcal I_Q$ 上提供信息；由此得到固定 $n$ 的矩还需要 $F_2\ge0$ 与有限减除的色散关系（\cref{thm:dis-fixed-n}），
而 $n<n_0$ 的低阶矩还需要增长条件（\cref{prop:dis-low-moments}）。}
\label{fig:dis-euclid-segment}
\end{figure}
%%% END fig:dis-euclid-segment
